% ============================================================================
% II. RELATED WORK  (target length: about 0.5 column; three short paragraphs)
% R2: originality versus RAG / uncertainty / decision-theoretic work must be
% explicit. Each paragraph ends with one sentence that states the difference
% from this paper. The submitted memory-system paragraphs (MemGPT, LongMem,
% MemoryBank, LightMem) are removed: Experiment 4 (memory) was cut and the
% 15-reference cap leaves no room for work the paper does not use.
% ============================================================================
\section{Related Work}\label{sec:related}

% ---- P-II.1 Adaptive retrieval ------------------------------------------------
% PURPOSE: position against work that decides WHETHER/HOW MUCH to retrieve.
% CLAIMS + EVIDENCE:
%   (a) RAG conditions generation on retrieved passages [lewis2020rag].
%   (b) FLARE triggers retrieval on low-confidence tokens of the upcoming
%       sentence [jiang2023flare]; DRAGIN on real-time information need
%       during generation [su2024dragin]; Self-RAG uses learned reflection
%       tokens [asai2024selfrag]; Adaptive-RAG routes by predicted query
%       complexity [jeong2024adaptiverag].
%   (c) Retrieval can lower accuracy on popular entities the model already
%       knows; popularity-aware adaptive retrieval [mallen2023trust]. This
%       motivates g^- > 0 in Eq.~(\ref{eq:obj}).
%   (d) DIFFERENCE: these methods decide whether to retrieve during text
%       generation; none prices the dispatch time against a user's ongoing
%       speech, and none asks when joint control of the two decisions adds
%       value over separate components.
% NOTE: the paper's B4 and B6 are single-turn, pre-generation proxies
%   "in the style of" FLARE and DRAGIN (thresholds on a confidence/entropy
%   statistic of a closed-book draft), not reimplementations; say so here or
%   in Sec. V [controller/baselines/b4_flare.py, b6_dragin.py].
\todo{P-II.1: adaptive retrieval \cite{lewis2020rag,jiang2023flare,su2024dragin,asai2024selfrag,jeong2024adaptiverag,mallen2023trust}
and the difference from this paper.}

% ---- P-II.2 Turn-taking ---------------------------------------------------------
% PURPOSE: position against work that predicts WHEN the user stops.
% CLAIMS + EVIDENCE:
%   (a) TurnGPT predicts turn shifts from text [ekstedt2020turngpt]; VAP
%       projects future voice activity from audio [ekstedt2022vap]; review of
%       end-of-turn detection and prediction [skantze2021turntaking].
%       (R1 minor: cite TurnGPT/VAP directly, not only the review.)
%   (b) Streaming ASR supplies the partial transcripts the belief is computed
%       on [radford2023whisper] -- only if v2 uses a Whisper-family model;
%       confirm against datasets/v2/*/speech/traces/*.jsonl asr_model field.
%   (c) DIFFERENCE: these predict the timing state Z (remaining speech); this
%       paper uses such a prediction as an input and asks how it should be
%       combined with the retrieval decision.
\todo{P-II.2: turn-taking and streaming ASR \cite{ekstedt2020turngpt,ekstedt2022vap,skantze2021turntaking,radford2023whisper}
and the difference from this paper.}

% ---- P-II.3 Decision theory -----------------------------------------------------
% PURPOSE: position against decision-theoretic control.
% CLAIMS + EVIDENCE:
%   (a) Value of information [howard1966information] (R1 minor) -- the
%       one-step waiting value I_k in Prop. 4 is Howard's VoI.
%   (b) Decision-theoretic turn-taking that weighs acting early against
%       waiting under uncertainty and processing delay [bohus2011decisions].
%   (c) Calibration and ECE [guo2017calibration] (used in Sec. VI).
%   (d) DIFFERENCE: the contribution is not the use of expected utility but
%       an exact characterisation of when a joint argmax and a decomposed
%       controller with shared prices coincide, and a test of the boundary.
\todo{P-II.3: value of information and decision-theoretic dialogue control
\cite{howard1966information,bohus2011decisions,guo2017calibration} and the
difference from this paper.}
